\documentclass[tikz,border=14pt]{standalone}
\usepackage[utf8]{inputenc}
\usepackage[T1]{fontenc}
\usepackage[french]{babel}
\usepackage{lmodern}
\usepackage{textcomp}
\usepackage{amssymb}
\usepackage{xcolor}
\usetikzlibrary{positioning,arrows.meta,shapes.geometric,calc,fit}

\definecolor{mainblue}{RGB}{28,74,102}
\definecolor{darkgray}{RGB}{62,62,66}
\definecolor{ent}{RGB}{221,236,247}
\definecolor{tab}{RGB}{215,237,221}
\definecolor{mod}{RGB}{233,233,238}
\definecolor{ext}{RGB}{250,233,207}
\definecolor{alerte}{RGB}{250,226,226}
\definecolor{kored}{RGB}{176,42,42}
\definecolor{okgreen}{RGB}{46,125,50}
\definecolor{nunfond}{RGB}{253,244,205}
\definecolor{nuntrait}{RGB}{176,128,0}
\definecolor{ndeuxfond}{RGB}{236,228,247}
\definecolor{ndeuxtrait}{RGB}{96,62,150}
% Interface : thème « Bordeaux & lin » de Timso AI
\definecolor{encre}{HTML}{241C21}
\definecolor{gris}{HTML}{6E6469}
\definecolor{trait}{HTML}{E8E1DE}
\definecolor{sable}{HTML}{FBF8F6}
\definecolor{bordeaux}{HTML}{7A2E3F}
\definecolor{bordeauxclair}{HTML}{F7ECEE}
\definecolor{vert}{HTML}{3A6B4C}
\definecolor{ambre}{HTML}{9C5510}

\newcommand{\f}[1]{\texttt{#1}}

\tikzset{
  boite/.style={rectangle,rounded corners=2pt,align=center,font=\scriptsize,inner sep=4pt,minimum height=0.75cm},
  etage/.style={boite,draw=mainblue,fill=ent},
  donnee/.style={boite,draw=mainblue!70,fill=tab},
  module/.style={boite,draw=darkgray!70,fill=mod},
  acteur/.style={boite,draw=darkgray,fill=ext},
  niv0/.style={boite,draw=mainblue,fill=ent,line width=0.9pt},
  niv1/.style={boite,draw=nuntrait,fill=nunfond,line width=0.9pt},
  niv2/.style={boite,draw=ndeuxtrait,fill=ndeuxfond,line width=0.9pt},
  repli/.style={boite,draw=kored!80,fill=alerte,dashed},
  test/.style={diamond,aspect=2.3,draw=mainblue,fill=white,align=center,font=\scriptsize,inner sep=1pt},
  bulleclient/.style={boite,draw=darkgray,fill=ext,text width=5.6cm,align=left},
  bulleassistant/.style={boite,draw=mainblue,fill=white,text width=5.6cm,align=left},
  json/.style={boite,draw=darkgray!60,fill=white,align=left,font=\ttfamily\scriptsize},
  fl/.style={-{Latex[length=1.8mm]},semithick,draw=mainblue!80},
  fl1/.style={-{Latex[length=1.8mm]},semithick,draw=nuntrait},
  fl2/.style={-{Latex[length=1.8mm]},semithick,draw=ndeuxtrait},
  flp/.style={-{Latex[length=1.8mm]},semithick,draw=darkgray!75,dashed},
  flr/.style={-{Latex[length=1.8mm]},semithick,draw=kored!75,dashed},
  lien/.style={semithick,draw=darkgray!60},
  et/.style={font=\scriptsize,text=darkgray,inner sep=1.5pt},
  titre/.style={font=\large\bfseries,text=mainblue},
  soustitre/.style={font=\scriptsize,text=darkgray},
  % interface
  carte/.style={draw=trait,fill=white,rounded corners=4pt},
  pastille/.style={draw=#1,fill=#1!9,text=#1,rounded corners=5pt,font=\tiny\bfseries,inner sep=2.5pt},
}

\begin{document}

% ════════════════════════════════════════════════════════════════════
% Page 1 — L'aiguilleur, vue d'ensemble
% ════════════════════════════════════════════════════════════════════
\begin{tikzpicture}[x=1cm,y=1cm]
  \node[titre] (t) at (0,1.5) {L'aiguilleur : le plus petit moyen qui suffit, message par message};
  \node[soustitre] at (0,1.0) {Vue d'ensemble --- chaque niveau est détaillé dans les pages suivantes};
  \node[acteur,text width=6cm] (msg) at (0,0) {\textbf{Message du client}\\ dans l'interface (vue démo ou vue utilisateur)};
  \node[module,text width=10.4cm] (entree) at (0,-1.85)
    {\textbf{Entrée commune, sans LLM} (environ 6 ms)\\[1pt]
     barrières : demande du code (sécurité), rendez-vous, relevés, au revoir, montant ambigu\\
     entités par regex (montants, durées, liens familiaux) · contexte : parcours en cours,\\
     question en attente, demandes en file, parcours mis de côté\\
     mots-clés (tous les sujets touchés) · e5-small + 2 régressions logistiques (confiance)};
  \node[module,text width=5.6cm,align=left] (signaux) at (9.6,-1.85)
    {\textbf{Les 5 signaux regardés par l'aiguilleur}\\[1pt]
     1. plusieurs demandes dans le message\\
     2. confiance sous $\tau = 0{,}76$\\
     3. réponse attendue illisible par les règles\\
     4. renvoi au contexte (\og mais ma carte ? \fg)\\
     5. question ouverte sur un sujet de nos fiches\\ \hphantom{5. }(\og pourquoi \fg, \og je peux ? \fg, \og plafond \fg)};
  \node[test] (aig) at (0,-4.2) {\textbf{Aiguilleur}\\ le plus petit moyen\\ qui suffit};
  \draw[lien,dashed] (signaux.south) |- (aig.east);
  \draw[fl] (msg) -- (entree);
  \draw[fl] (entree) -- (aig);

  \node[niv0,text width=4.8cm] (n0) at (-6.2,-6.7)
    {\textbf{Niveau 0 --- Règles}\\ demande claire, comprise avec confiance,\\ ou réponse lue par les règles\\
     \textbf{76\,\% des messages, 0 LLM}};
  \node[niv1,text width=4.8cm] (n1) at (0,-6.7)
    {\textbf{Niveau 1 --- Petit LLM interprète}\\ signaux 2 à 4 : avec \textbf{8 exemples proches}\\ (RAG d'exemples), il rend un \textbf{JSON}\\
     \textbf{imposé par une grammaire} ; il ne\\ \textbf{répond jamais au client}};
  \node[niv2,text width=4.8cm] (n2) at (6.2,-6.7)
    {\textbf{Niveau 2 --- RAG + LLM rédacteur}\\ signal 5 : question ouverte sur un\\ produit ou une règle ; une fois, à la fin,\\
     \textbf{uniquement depuis nos 37 fiches}};
  \draw[fl] (aig.west) -| node[et,above,pos=0.3]{clair} (n0.north);
  \draw[fl1] (aig.south) -- node[et,right]{flou} (n1.north);
  \draw[fl2] (aig.east) -- ++(1.6,0) -| node[et,above,pos=0.2]{question ouverte} (n2.north);

  \node[etage,text width=4.8cm] (auto) at (-6.2,-9.2)
    {\textbf{Automate + 9 parcours YAML}\\ file par urgence, mémoire, reprise après\\ une parenthèse ou un refus, corrections,\\ vraisemblance des montants ;\\ règles de conseil (HCSF, plafonds)};
  \node[etage,text width=4.8cm] (gab) at (-6.2,-11.1)
    {\textbf{Gabarits}\\ hypothèses du calcul, mention légale,\\ glossaire, mots simples};
  \draw[fl] (n0) -- (auto);
  \draw[fl] (auto) -- (gab);
  \draw[fl1] ($(n1.south)+(-1.4,0)$) |- node[et,below,pos=0.62]{retour au niveau 0} (auto.east);
  \node[repli,text width=3.8cm] (clar) at (1.4,-11.1)
    {rien d'exploitable, ou LLM\\ indisponible (10 s) : niveau 0,\\ \textbf{question de clarification}};
  \draw[flr] ($(n1.south)+(1.4,0)$) -- (clar.north);

  \node[niv2,text width=4.8cm] (rag) at (6.2,-9.2)
    {\textbf{Recherche dans nos 37 fiches}\\ e5-small, 3 passages les plus proches\\ (produits, règles, glossaire)};
  \node[test] (gf) at (6.2,-11.1) {garde-fou\\ réussi ?};
  \draw[fl2] (n2) -- (rag);
  \draw[fl2] (rag) -- node[et,right]{le LLM rédige} (gf);
  \node[repli,text width=3.6cm] (rien) at (11.0,-9.2)
    {cosinus $<$ 0,82 : \og je n'ai pas\\ cette information \fg{} + conseiller};
  \draw[flr] (rag) -- (rien);
  \node[repli,text width=3.6cm] (extrait) at (11.0,-11.1)
    {échec : extrait de la meilleure\\ fiche + sources ; le texte\\ écarté reste consultable};
  \draw[flr] (gf) -- node[et,above]{non} (extrait);

  \node[acteur,text width=14.2cm] (out) at (0,-13.2)
    {\textbf{Réponse au client}\\ + niveau et raison, modèle, latence, énergie, CO$_2$ estimé
     + volet \og Pourquoi \fg{} (intention, confiance, règles, rédaction écartée) + volet \og Sources \fg{} (passages, scores)};
  \draw[fl] (gab.south) -- (gab.south |- out.north);
  \draw[flr] (clar.south) -- (clar.south |- out.north);
  \draw[fl2] (gf.south) -- node[et,right]{oui} (gf.south |- out.north);

  \node[font=\scriptsize,text=darkgray,align=center] at (0,-14.5)
    {\textcolor{mainblue}{\rule{6pt}{6pt}} niveau 0, règles \quad
     \textcolor{nuntrait}{\rule{6pt}{6pt}} niveau 1, petit LLM qui comprend \quad
     \textcolor{ndeuxtrait}{\rule{6pt}{6pt}} niveau 2, RAG et LLM qui rédige \quad
     \textcolor{kored}{\rule{6pt}{6pt}} replis sûrs : jamais de devinette\\[2pt]
     Un seul modèle pour les niveaux 1 et 2 : \f{qwen3:1.7b} (1,4 Go). La décision est prise à chaque message :
     une conversation passe de 1 à 2, puis à 0.};
\end{tikzpicture}

% ════════════════════════════════════════════════════════════════════
% Page 2 — Niveau 0, sur le début du scénario de démonstration
% ════════════════════════════════════════════════════════════════════
\begin{tikzpicture}[node distance=0.32cm]
  \node[bulleclient] (c1) {\textbf{Client} : J'ai perdu ma carte et je veux faire un virement.};
  \node[niv0,text width=5.6cm,right=1.2cm of c1] (i1)
    {mots-clés : \textbf{2 sujets} (carte perdue, virement)\\ $\rightarrow$ file par urgence : 1. carte perdue, 2. virement};
  \node[bulleassistant,below=of c1] (a1)
    {\textbf{Assistant} : Je comprends, il faut agir vite. Faites opposition tout de suite : dans votre
     espace bancaire sécurisé, rubrique Opposition, ou en appelant votre conseiller. [\dots]
     Rappel : la banque ne vous demandera jamais votre code confidentiel.\\
     Passons maintenant à votre autre demande : un virement. [\dots] À qui voulez-vous envoyer l'argent ?};
  \node[niv0,text width=5.6cm] (i2) at (i1 |- a1)
    {parcours \og carte perdue \fg{} : consigne \textbf{immédiate},\\ aucun accès au compte (périmètre du kick-off),\\
     puis la demande suivante de la file : le virement};
  \node[bulleclient,below=of a1] (c2) {\textbf{Client} : c'est pour mon papa\\
     \textit{(le \og oui \fg{} qui le précède passe au niveau 1 : page suivante)}};
  \node[niv0,text width=5.6cm] (i3) at (i1 |- c2)
    {réponse attendue lue par les règles :\\ \og papa \fg{} $\rightarrow$ bénéficiaire = père (mémorisé)};
  \node[bulleassistant,below=of c2] (a2) {\textbf{Assistant} : Quel montant voulez-vous envoyer ?};
  \node[bulleclient,below=of a2] (c3) {\textbf{Client} : 100000000};
  \node[niv0,text width=5.6cm] (i4) at (i1 |- c3)
    {\textbf{vraisemblance} : au moins 10 fois les revenus\\ mensuels $\rightarrow$ alerte, avec le solde du profil (1\,450\,\texteuro)};
  \node[bulleassistant,below=of c3] (a3)
    {\textbf{Assistant} : Attention : 100\,000\,000\,\texteuro, c'est un montant très inhabituel, bien au-delà
     de votre solde (1\,450\,\texteuro) : vérifiez-le avant de valider. Pour effectuer ce virement, utilisez votre
     espace bancaire sécurisé [\dots] Je ne peux pas accéder à votre compte ni exécuter le virement.};
  \node[niv0,text width=5.6cm] (i5) at (i1 |- a3)
    {parcours \og virement \fg{} terminé : orientation\\ vers l'espace sécurisé, sans accès au compte};
  \foreach \a/\b in {c1/i1,a1/i2,c2/i3,c3/i4,a3/i5} \draw[lien,dashed] (\a.east) -- (\b.west);

  \node[module,text width=6.7cm,align=left,right=0.9cm of i3] (outils)
    {\textbf{Ce qui travaille au niveau 0}\\[2pt]
     · barrières : demande du code, rendez-vous, relevés\\
     · automate de dialogue et 9 parcours YAML\\
     · file des demandes, par ordre d'urgence\\
     · mémoire : reprise après une parenthèse\\
     · reprise après un refus (\og finalement oui \fg)\\
     · corrections (\og plutôt 8000 \fg) : conseil recalculé\\
     · vraisemblance des montants\\
     · règles de conseil (HCSF, plafonds)\\
     · gabarits, glossaire, mots simples\\[2pt]
     \textbf{Coût} : 0 LLM, quelques ms ;\\
     0,4 Wh pour 1\,000 messages (CodeCarbon)};

  \node[titre,above=0.9cm of c1.north east,anchor=south] (t) {Niveau 0 --- Règles : la majorité des messages, sans LLM};
  \node[soustitre,below=0.12cm of t] {Le début du scénario de démonstration (Inès, C01) --- à gauche la conversation, à droite ce qui se passe};
\end{tikzpicture}

% ════════════════════════════════════════════════════════════════════
% Page 3 — Niveau 1 : le petit LLM qui interprète
% ════════════════════════════════════════════════════════════════════
\begin{tikzpicture}[node distance=0.42cm]
  \node[module,text width=8.4cm,align=left] (ctx)
    {\textbf{Contexte transmis au LLM}, en trois lignes --- \textbf{jamais le profil}\\[1pt]
     sujet en cours : un virement\\
     question en attente : \og À qui voulez-vous envoyer l'argent ? \fg\\
     déjà traité : une carte perdue ou volée};
  \node[bulleclient,text width=7.0cm,below=of ctx] (m)
    {\textbf{Client} : mais ma carte ?\\
     \textit{(ou : \og oui \fg, \og perdue, mais je veux aussi faire un virement \fg)}};
  \node[niv1,text width=7.2cm,below=of m] (pourquoi)
    {\textbf{Pourquoi le niveau 1} : renvoi au contexte, confiance sous $\tau$,\\
     ou réponse attendue que les règles ne savent pas lire};
  \node[niv1,text width=7.2cm,below=of pourquoi] (ex)
    {\textbf{RAG d'exemples} : e5-small cherche les \textbf{8 phrases d'entraînement}\\
     \textbf{les plus proches} et les donne au LLM avec leur sujet};
  \node[json,text width=5.6cm,right=0.8cm of ex] (exliste)
    {« Ma carte ne fonctionne plus » → carte\_usage\\
     « Ma carte a peut-être été piratée » → carte\_perdue\\
     « Est-ce que ma carte marche au Maroc ? » →\\ \ \ change\_etranger \quad [\dots] (8 au total)};
  \draw[lien,dashed] (exliste) -- (ex);
  \node[niv1,text width=7.2cm,below=of ex] (llm)
    {\textbf{\f{qwen3:1.7b}} (1,4 Go), un seul appel, délai maximum 10 s\\
     sortie \textbf{imposée par une grammaire} (schéma JSON : Ollama \f{format},\\ vLLM \f{response\_format}) : 100\,\% de JSON valide};
  \node[json,text width=8.4cm,below=of llm] (j)
    {\{ "intentions": ["carte\_usage"],\\
     \ \ "reponse\_a\_la\_question": "ma carte ?" \}\\[2pt]
     \textrm{\textit{1 à 3 sujets d'une liste fermée : nos 8 intentions et les familles}}\\
     \textrm{\textit{de Banking77 (ou \og aucune \fg) ; le code les traduit en intentions}}};
  \node[test,below=0.5cm of j] (v) {vérifié par\\ le code ?};
  \node[module,text width=6.6cm,align=left,right=0.8cm of j] (verifs)
    {\textbf{Ce que vérifie le code}\\[1pt]
     · sujets de la liste fermée, réordonnés par urgence\\
     · la réponse lue reprend les mots du message\\
     · aucun nombre inventé\\
     · renvoi \og ma carte \fg{} : rattaché à la carte perdue\\
     \ \ déjà traitée\\
     · nouvelle demande pendant une question :\\
     \ \ confirmée par les mots-clés ou le classifieur};
  \draw[lien,dashed] (verifs.south) |- (v.east);
  \node[niv0,text width=5.0cm] (ok) at ($(v.south)+(-3.6,-1.5)$)
    {\textbf{Retour au niveau 0}\\ l'automate range la compréhension\\ (file, mémoire), les gabarits répondent :\\ ici, la consigne d'opposition};
  \node[repli,text width=4.4cm] (ko) at ($(v.south)+(3.6,-1.5)$)
    {\textbf{Rien d'exploitable} : niveau 0,\\ question de clarification\\ \og votre carte, ou un virement ? \fg};
  \draw[fl1] (ctx) -- (m);
  \draw[fl1] (m) -- (pourquoi);
  \draw[fl1] (pourquoi) -- (ex);
  \draw[fl1] (ex) -- (llm);
  \draw[fl1] (llm) -- (j);
  \draw[fl1] (j) -- (v);
  \draw[fl] (v.west) -| node[et,above,pos=0.2]{oui} (ok.north);
  \draw[flr] (v.south) -- ++(0,-0.4) -| node[et,above,pos=0.3]{non} (ko.north);

  \node[module,text width=4.6cm,align=left,left=0.8cm of llm] (regles)
    {\textbf{Les règles du niveau 1}\\[2pt]
     · le LLM ne voit \textbf{pas le profil}\\
     · il ne \textbf{répond jamais} au client\\
     · il n'écrit \textbf{aucun chiffre} que le\\ \hphantom{· }client n'a pas écrit\\
     · s'il est indisponible : niveau 0\\[2pt]
     \textbf{Mesuré} (36 messages difficiles) :\\ 30 justes avec le RAG d'exemples,\\ 15 sans ; p50 188 ms (DeepEirb)};

  \node[titre,above=0.9cm of ctx] (t) {Niveau 1 --- Un petit LLM qui comprend, sans répondre};
  \node[soustitre,below=0.12cm of t] {Il remplace la devinette par une compréhension structurée, que le code vérifie avant de s'en servir};
\end{tikzpicture}

% ════════════════════════════════════════════════════════════════════
% Page 4 — Niveau 2 : RAG + LLM rédacteur
% ════════════════════════════════════════════════════════════════════
\begin{tikzpicture}[x=1cm,y=1cm]
  \node[titre] (t) at (6.2,2.2) {Niveau 2 --- RAG : chercher dans nos fiches, puis rédiger en citant};
  \node[soustitre] at (6.2,1.7) {Le même encodeur que pour comprendre ; un LLM seulement quand il faut rédiger, et sous garde-fou};
  % Au démarrage
  \node[donnee,text width=3.5cm] (fiches) at (0,0)
    {\textbf{37 fiches sourcées}\\ 11 produits du catalogue,\\ 4 règles (HCSF, notaire, frais),\\ 18 termes du glossaire,\\ 4 fiches écrites à la main};
  \node[niv2,text width=2.9cm] (decoupe) at (4.1,0) {\textbf{Une fiche,}\\ \textbf{un passage}\\ avec sa source\\ et sa date};
  \node[niv2,text width=2.9cm] (encode) at (7.9,0) {\textbf{e5-small}\\ préfixe \f{passage:}\\ (déjà chargé : aucun\\ modèle en plus)};
  \node[niv2,text width=2.9cm] (index) at (11.7,0) {\textbf{Index en mémoire}\\ 37 vecteurs,\\ sans base vectorielle};
  \draw[fl2] (fiches) -- (decoupe);
  \draw[fl2] (decoupe) -- (encode);
  \draw[fl2] (encode) -- (index);
  \node[draw=darkgray!60,dashed,rounded corners=3pt,fit=(fiches)(index),inner sep=7pt] (cadre) {};
  \node[et,anchor=south west] at (cadre.north west) {\textit{Au démarrage de l'application, une fois (environ 1 s)}};

  % À chaque question ouverte
  \node[bulleclient,text width=5.2cm] (q) at (2,-3.0)
    {\textbf{Client} : Je peux retirer de mon PEL avant 4 ans ?};
  \node[niv2,text width=5.0cm] (qe) at (2,-4.45) {\textbf{e5-small}, préfixe \f{query:}, similarité\\ avec chaque passage ; \textbf{recherche hybride} :\\ bonus si la question nomme le sujet};
  \node[niv2,text width=5.0cm] (top) at (2,-5.75) {\textbf{3 passages les plus proches}, avec leur score};
  \node[test] (seuil) at (2,-7.15) {cosinus\\ $\geq$ 0,82 ?};
  \node[niv2,text width=5.0cm] (red) at (2,-8.8)
    {\textbf{\f{qwen3:1.7b} rédige}, une fois, délai max 20 s\\ consigne : répondre \textbf{uniquement} avec\\ les passages, les citer [1] [2], rester bref};
  \node[test] (gf) at (2,-10.45) {garde-fou\\ réussi ?};
  \draw[fl2] (q) -- (qe);
  \draw[fl2] (qe) -- (top);
  \draw[fl2] (top) -- (seuil);
  \draw[fl2] (seuil) -- node[et,right]{oui} (red);
  \draw[fl2] (red) -- (gf);
  \draw[flp] (index.south) |- node[et,above,pos=0.75]{comparaison} (qe.east);

  \node[repli,text width=3.6cm] (rien) at (6.9,-7.15)
    {\og Je n'ai pas cette information\\ dans mes fiches \fg{} + conseiller\\ --- jamais d'invention};
  \draw[flr] (seuil) -- node[et,above]{non} (rien);
  \node[module,text width=6.0cm,align=left] (verifs) at (8.9,-10.45)
    {\textbf{Le garde-fou}\\[1pt]
     · \textbf{tous les chiffres} viennent des passages\\
     · \textbf{fidélité} : au moins 80\,\% des mots de la\\
     \ \ réponse viennent des passages\\
     · aucune formule interdite, longueur plausible\\
     · une citation (ajoutée si la réponse est fidèle)};
  \draw[lien,dashed] (verifs) -- (gf);
  \node[module,text width=5.2cm,align=left] (quand) at (12.7,-7.3)
    {\textbf{Quand le niveau 2 ?}\\[1pt]
     · question ouverte sur un produit\\
     \ \ ou une règle, hors parcours\\
     · \og pourquoi ? \fg{} après un conseil : les\\
     \ \ passages sont les phrases du conseil\\
     \ \ (\og Votre simulation \fg)\\[2pt]
     \textbf{Les chiffres} restent ceux des fiches\\
     et des règles : le LLM ne calcule rien};

  \node[bulleassistant,text width=6.4cm] (rep) at (-0.6,-13.5)
    {\textbf{Assistant} : Oui, mais les fonds ne sont pas définitivement bloqués. Le retrait anticipé
     entraîne la clôture et des conséquences selon l'ancienneté [1].\\
     Revenons à votre virement : quel montant voulez-vous envoyer ?\\[2pt]
     \textit{Sources : [1] Plan épargne logement (Service-Public, 27/10/2026) ; [2] PEL (glossaire) ; [3] LEP}};
  \node[repli,text width=4.6cm] (ext) at (6.6,-13.5)
    {\textbf{Repli} : \og D'après nos fiches : \fg{} l'extrait\\ de la meilleure fiche, avec ses sources,\\ sans LLM. Le texte écarté est montré\\ dans \og Pourquoi \fg, mots absents surlignés};
  \draw[fl2] (gf.west) -| node[et,above,pos=0.2]{oui} (rep.north);
  \draw[flr] (gf.east) -- ++(0.35,0) |- ($(ext.north)+(0,0.35)$) -- node[et,right,pos=0.1]{non} (ext.north);
  \node[font=\scriptsize,text=darkgray,align=left,anchor=north west] at (9.6,-12.7)
    {\textbf{Mesuré} (22 questions, DeepEirb) :\\ bonne fiche dans les 3 passages : 18/18\\ hors de nos fiches, refusé : 4/4\\ réponses justes : 22/22 ; p50 391 ms};
\end{tikzpicture}

% ════════════════════════════════════════════════════════════════════
% Page 5 — L'interface, dans le style de Timso AI
% ════════════════════════════════════════════════════════════════════
\begin{tikzpicture}[x=1cm,y=1cm,font=\scriptsize]
  \fill[white] (0,0) rectangle (17,11.2);
  \draw[trait] (0,0) rectangle (17,11.2);
  % En-tête
  \draw[trait] (0,10.4) -- (17,10.4);
  \node[anchor=west,font=\normalsize\bfseries,text=encre] at (0.3,10.8) {Assistant \textcolor{bordeaux}{bancaire}};
  \draw[trait,fill=sable,rounded corners=4pt] (5.0,10.55) rectangle (10.9,11.05);
  \fill[bordeauxclair,rounded corners=3pt] (5.08,10.62) rectangle (7.6,10.98);
  \node[text=bordeaux,font=\tiny\bfseries] at (6.34,10.8) {Démo : tout voir};
  \node[text=gris,font=\tiny] at (9.25,10.8) {Utilisateur : les messages};
  \fill[bordeauxclair,rounded corners=3pt] (11.1,10.62) rectangle (12.85,10.98);
  \node[text=bordeaux,font=\tiny\bfseries] at (11.97,10.8) {Scénarios};
  \node[text=gris,font=\tiny] at (13.6,10.8) {Mode sombre};
  \draw[trait,rounded corners=3pt] (14.5,10.58) rectangle (16.9,11.02);
  \node[text=encre,font=\tiny\bfseries] at (15.7,10.8) {Nouvelle conversation};

  % Colonne de réglages
  \draw[carte] (0.3,7.0) rectangle (4.2,10.1);
  \node[anchor=north west,text=gris,font=\tiny\bfseries] at (0.45,10.0) {PROFIL SIMULÉ};
  \draw[trait,rounded corners=3pt] (0.45,9.25) rectangle (4.05,9.65);
  \node[anchor=west,text=encre,font=\tiny] at (0.55,9.45) {C01 --- Inès, 34 ans --- cadre};
  \node[anchor=north west,text=encre,align=left,font=\tiny] at (0.45,9.1)
    {\textbf{Inès}\\ \textcolor{gris}{ÂGE} \ 34\\ \textcolor{gris}{REVENUS NETS} \ 3\,200 \textcolor{ambre}{simulé}\\ \textcolor{gris}{ÉPARGNE} \ 8\,000 \textcolor{ambre}{simulé}\\ \textcolor{gris}{SOLDE} \ 1\,450 \textcolor{ambre}{simulé}};
  \draw[carte] (0.3,5.2) rectangle (4.2,6.8);
  \node[anchor=north west,text=gris,font=\tiny\bfseries] at (0.45,6.7) {AIGUILLAGE};
  \draw[trait,rounded corners=3pt] (0.45,5.95) rectangle (4.05,6.35);
  \node[anchor=west,text=encre,font=\tiny] at (0.55,6.15) {Automatique};
  \node[anchor=west,text=encre,font=\tiny] at (0.45,5.55) {\textcolor{bordeaux}{$\blacksquare$} LLM actif \quad $\square$ Mots simples};
  \draw[carte] (0.3,2.5) rectangle (4.2,5.0);
  \node[anchor=north west,text=gris,font=\tiny\bfseries] at (0.45,4.9) {CETTE CONVERSATION};
  \node[anchor=north west,text=encre,font=\tiny\bfseries] at (0.45,4.55) {4 message(s)};
  \node[pastille=vert] at (0.75,3.95) {0}; \node[text=gris,font=\tiny] at (1.0,3.95) {→};
  \node[pastille=ambre] at (1.25,3.95) {1}; \node[text=gris,font=\tiny] at (1.5,3.95) {→};
  \node[pastille=vert] at (1.75,3.95) {0}; \node[text=gris,font=\tiny] at (2.0,3.95) {→};
  \node[pastille=bordeaux] at (2.25,3.95) {2};
  \node[anchor=north west,text=gris,align=left,font=\tiny] at (0.45,3.7)
    {appels au LLM : 2\\ latence : 0,6 s\\ énergie estimée : 46 mWh\\ \textcolor{vert}{carbone estimé : 2,6 mg CO$_2$e}};
  \draw[carte] (0.3,0.3) rectangle (4.2,2.4);
  \node[anchor=north west,text=gris,font=\tiny\bfseries] at (0.45,2.3) {EXEMPLES, SYSTÈME};
  \node[anchor=north west,text=gris,align=left,font=\tiny] at (0.45,1.95)
    {LLM : \textcolor{vert}{serveur joignable}\\ niveaux 1 et 2 : \texttt{qwen3:1.7b}\\ classifieur : e5-small, $\tau$ = 0,76\\ RAG : 37 fiches sourcées};

  % Conversation
  \fill[bordeaux,rounded corners=6pt] (11.0,9.55) rectangle (16.7,10.1);
  \node[text=white,anchor=east,font=\tiny] at (16.5,9.83) {Je peux retirer de mon PEL avant 4 ans ?};
  \fill[sable,rounded corners=3pt] (4.6,5.0) rectangle (14.9,9.3);
  \draw[trait,rounded corners=3pt] (4.6,5.0) rectangle (14.9,9.3);
  \fill[bordeaux] (4.6,5.0) rectangle (4.68,9.3);
  \node[text=encre,anchor=north west,align=left,text width=9.8cm,font=\tiny] at (4.85,9.15)
    {Oui, mais les fonds ne sont pas définitivement bloqués. Le retrait anticipé entraîne la clôture et des
     conséquences selon l'ancienneté. \textcolor{bordeaux}{[1]}\\[3pt] Revenons à votre virement : quel montant voulez-vous envoyer ?};
  \fill[white] (4.68,7.55) rectangle (14.9,8.05);
  \draw[trait] (4.68,8.05) -- (14.9,8.05);
  \node[pastille=bordeaux,anchor=west] at (4.85,7.8) {Niveau 2 · RAG et rédaction};
  \node[anchor=west,text=gris,font=\tiny] at (8.75,7.8) {\texttt{qwen3:1.7b} \quad 391 ms \quad $\approx$ 22,9 mWh};
  \node[anchor=west,text=vert,font=\tiny\bfseries] at (12.7,7.8) {$\approx$ 1,3 mg CO$_2$e};
  \fill[white] (4.68,5.0) rectangle (14.9,7.55);
  \draw[trait] (4.68,7.55) -- (14.9,7.55);
  \node[anchor=west,text=gris,font=\tiny] at (4.85,7.3) {\textcolor{bordeaux}{+} Pourquoi};
  \draw[trait] (4.68,7.05) -- (14.9,7.05);
  \node[anchor=west,text=gris,font=\tiny] at (4.85,6.8) {\textcolor{bordeaux}{--} Sources (3)};
  \foreach \y/\n/\titre/\s in {6.3/1/{PLAN ÉPARGNE LOGEMENT (PEL)}/{0,904}, 5.75/2/{PEL (GLOSSAIRE)}/{0,870}, 5.25/3/{LIVRET D'ÉPARGNE POPULAIRE (LEP)}/{0,847}} {
    \draw[trait] (4.68,\y+0.25) -- (14.9,\y+0.25);
    \node[anchor=west,text=gris,font=\tiny\bfseries] at (4.85,\y) {PASSAGE \n\ · \titre};
    \node[anchor=east,text=bordeaux,font=\tiny\bfseries] at (14.75,\y) {\s};
  }
  \node[anchor=west,text=gris,font=\tiny,align=left] at (15.15,7.8) {chaque\\ réponse :\\ niveau,\\ énergie,\\ CO$_2$};

  % Barre de scénario
  \fill[sable,rounded corners=3pt] (4.6,1.25) rectangle (16.7,2.55);
  \draw[trait,rounded corners=3pt] (4.6,1.25) rectangle (16.7,2.55);
  \fill[bordeaux] (4.6,1.25) rectangle (4.68,2.55);
  \node[anchor=west,text=bordeaux,font=\tiny\bfseries] at (4.85,2.25) {SCÉNARIO};
  \draw[trait,fill=white,rounded corners=3pt] (6.5,2.05) rectangle (10.4,2.45);
  \fill[bordeauxclair,rounded corners=2pt] (6.55,2.09) rectangle (7.75,2.41);
  \node[text=bordeaux,font=\tiny\bfseries] at (7.15,2.25) {1 · Inès};
  \node[text=gris,font=\tiny] at (8.4,2.25) {2 · Léa};
  \node[text=gris,font=\tiny] at (9.6,2.25) {3 · Karim};
  \foreach \i in {0,...,9} {
    \pgfmathsetmacro{\x}{13.6+\i*0.2}
    \ifnum\i<3 \fill[vert] (\x,2.25) circle (0.06); \else \ifnum\i=3 \fill[bordeaux] (\x,2.25) circle (0.06); \else \fill[trait] (\x,2.25) circle (0.06); \fi\fi
  }
  \node[anchor=west,text=gris,font=\tiny] at (15.6,2.25) {3 / 10};
  \draw[trait,dashed,fill=white,rounded corners=3pt] (4.85,1.4) rectangle (14.9,1.85);
  \node[anchor=west,text=encre,font=\tiny] at (4.95,1.62) {\textcolor{gris}{\textbf{PROCHAIN}} \ « Je peux retirer de mon PEL avant 4 ans ? »};
  \fill[bordeaux,rounded corners=3pt] (15.1,1.4) rectangle (16.55,1.85);
  \node[text=white,font=\tiny\bfseries] at (15.82,1.62) {Envoyer ›};
  % Saisie
  \draw[trait] (4.4,1.05) -- (16.9,1.05);
  \draw[trait,rounded corners=4pt,fill=white] (4.6,0.3) rectangle (15.0,0.85);
  \node[anchor=west,text=gris,font=\tiny] at (4.8,0.57) {Votre message\dots};
  \fill[bordeaux,rounded corners=4pt] (15.15,0.3) rectangle (16.7,0.85);
  \node[text=white,font=\tiny\bfseries] at (15.92,0.57) {Envoyer};

  \node[titre,anchor=south] at (8.5,11.85) {L'interface de la démo, dans le style de Timso AI};
  \node[soustitre,anchor=south,align=center,text width=16cm] at (8.5,11.35)
    {Page web (FastAPI + HTML), thème \og Bordeaux \& lin \fg{} : vue démo (tout voir) ou vue utilisateur (les messages seuls) ;
     chaque réponse montre son niveau, sa raison, son énergie, son CO$_2$ et ses sources ; la barre de scénario joue la démo d'un clic};
\end{tikzpicture}

\end{document}
